\chapter{Spheres, Hyperbolic Space, and the Meaning of \(k\)}

Let us do a little mathematics, some simple mathematics. Spheres in any number of dimensions are special geometries: they are uniform, and here uniform means constant curvature. Positive curvature comes first, then uniform negative curvature, and after that the connection with Einstein's equations.

A spherical universe means that \emph{space} is spherical, not spacetime. There is also the other possibility, uniform negative curvature. Putting these cases side by side is what makes the role of general relativity clear.

\section{Spheres and Uniform Geometry}

A one-sphere is just a circle. The simplest way to describe it is to embed it in a two-dimensional auxiliary space and write
\begin{equation}
x^2+y^2=R^2.
\end{equation}
The unit one-sphere is
\begin{equation}
x^2+y^2=1.
\end{equation}

A two-sphere is described the same way, now using three auxiliary coordinates:
\begin{equation}
x^2+y^2+z^2=R^2,
\end{equation}
with unit case
\begin{equation}
x^2+y^2+z^2=1.
\end{equation}

These coordinates are only a crutch. The geometry of interest is not the surrounding \(x,y,z\) space. It is the surface itself. A creature living on a one-sphere lives on a circle; it has no access to an inside or an outside. The same idea applies to a two-sphere. The filling of the sphere and the outside of the sphere belong only to the auxiliary picture.

The same construction gives a three-sphere. One introduces one more coordinate, say \(w\), and writes
\begin{equation}
x^2+y^2+z^2+w^2=R^2,
\end{equation}
with unit case
\begin{equation}
x^2+y^2+z^2+w^2=1.
\end{equation}
A three-sphere is harder to visualize, because the crutch would now be four-dimensional, but the mathematics is not much harder.

The standard notation is \(\Omega_n\) for the \(n\)-sphere. Thus \(\Omega_1\) is the one-sphere, \(\Omega_2\) the two-sphere, and \(\Omega_3\) the three-sphere.

\section{Metrics on the Sphere}

The more useful description is intrinsic. One describes the sphere by coordinates that live on the sphere itself.

For the unit one-sphere one coordinate \(\phi\) is enough, with \(0\leq \phi \leq 2\pi\). Its metric is
\begin{equation}
ds^2=d\phi^2.
\end{equation}
This is also written
\begin{equation}
d\Omega_1^2=d\phi^2.
\end{equation}

The unit two-sphere is described as a sequence of one-spheres. One angle \(\phi\) still goes around each circle, but a second angle \(\theta\) runs from one end of the sphere to the other. That second angle goes from \(0\) to \(\pi\). Motion in the \(\theta\)-direction contributes \(d\theta^2\). Motion around the \(\phi\)-direction takes place on a circle whose radius is \(\sin\theta\). Thus
\begin{equation}
d\Omega_2^2=d\theta^2+\sin^2\theta\,d\phi^2
\end{equation}
or
\begin{equation}
d\Omega_2^2=d\theta^2+\sin^2\theta\,d\Omega_1^2.
\end{equation}

The three-sphere is obtained the same way, now as a sequence of two-spheres. Introduce another coordinate \(\alpha\). Then
\begin{equation}
d\Omega_3^2=d\alpha^2+\sin^2\alpha\,d\Omega_2^2.
\end{equation}
Expanded out, this is
\begin{equation}
d\Omega_3^2=d\alpha^2+\sin^2\alpha\left(d\theta^2+\sin^2\theta\,d\phi^2\right).
\end{equation}
Here \(\alpha\) runs from one zero-size end to the other, so \(0\leq \alpha \leq \pi\).\editorialnote{At 00:17:24 the audio briefly sounds like ``zero to two pi,'' but the surrounding construction and the explicit \(\sin^2\alpha\) factor require the new coordinate to run from one zero-size end of the family of lower-dimensional spheres to the other.}

The general recursive form is
\begin{equation}
d\Omega_n^2=d\alpha^2+\sin^2\alpha\,d\Omega_{n-1}^2.
\end{equation}

Once this point of view is adopted, the embedding picture can be forgotten. A one-sphere is described by one coordinate, a two-sphere by two, a three-sphere by three, and so on. The geometry stands on its own.

This geometry is special because it is homogeneous. One could deform the metric by replacing \(\sin^2\theta\) with some other function that still vanishes at both ends. One would then get some deformed closed surface, for example an ellipsoid. But it would no longer be uniform. The sphere is homogeneous because every point is the same as every other point.

If one wants a sphere of radius \(A\) instead of unit radius, one multiplies the unit metric by \(A^2\):
\begin{equation}
ds^2=A^2\,d\Omega_3^2.
\end{equation}
The square is necessary because the metric measures squared distance.

A brief aside concerns the zero-sphere. Extrapolating one step backward gives
\begin{equation}
x^2=1,
\end{equation}
whose solutions are the two points \(x=\pm 1\). So the unit zero-sphere is just two points.

\section{A Spherical FRW Universe}

Now the interest is not in space alone, but in spacetime. The relevant picture is a three-sphere whose radius changes with time.

One warning matters throughout: time is different. One should not think of the expanding universe as a higher-dimensional sphere with time playing the role of one more embedding coordinate.

With that in mind, the spacetime interval is
\begin{equation}
d\tau^2=dt^2-a(t)^2\,d\Omega_3^2.
\end{equation}
This is the FRW metric for a spatially spherical universe. The only new ingredient compared with the static sphere is that the radius has become a time-dependent function \(a(t)\).

It is convenient to choose coordinates so that we are at one pole. Nothing is special about that point; it is just a convention. Then \(\alpha\) can be thought of as a radial coordinate away from us on the unit sphere. It is not a radial distance in kilometers. It is a dimensionless angular measure on the unit sphere, but it plays the role of radial distance in the geometry.

Apart from local lumps such as galaxies and stars, this is the whole geometry of a spatially spherical FRW universe. Its spatial curvature is positive, and the sign convention is written
\begin{equation}
k=+1.
\end{equation}
Nothing itself is ``plus one.'' It simply means positive spatial curvature.

\section{Looking Out from a Sphere}

How does such a sphere look when one looks out from one point on it? This is where the stereographic projection enters.

Take a two-sphere tangent to an infinite plane at the south pole. Pick a point on the sphere. Draw the straight line from the north pole through that point and continue it until it hits the plane. In this way each point on the sphere is mapped to a point on the plane. The south pole lands at the origin of the plane. The north pole is the exceptional point: as one approaches it from different directions, its image runs off to infinity in a direction-sensitive way. One therefore treats it as the point at infinity.

If equal ``continents'' are placed at different positions on the sphere, their images on the plane become larger as one moves toward the north pole. So an observer placed at the south pole would see distant objects subtend larger angles than they would in flat space. In positive curvature, a known object looks too big relative to Euclidean geometry.

That is the basic observational point. Curvature is inferred from the relation between distance, size, and angular size. Galaxies themselves are not good standard rulers for very deep distances, but the logic is exactly that.

\begin{classroomqa}
\lecturetimestamp{00:45:54}
\audiencequestion{Do you also have to know the distance, not just the angular size?}
\lecturerresponse{Yes. The angle by itself is not enough. You need the angle and the distance.}
\end{classroomqa}

\begin{classroomqa}
\lecturetimestamp{00:46:18}
\audiencequestion{Do circles on the surface of the sphere project as circles?}
\lecturerresponse{Yes. Circles project to circles. That is one of the special features of stereographic projection.}
\end{classroomqa}

\begin{classroomqa}
\lecturetimestamp{00:47:02}
\audiencequestion{Would a continent keep the same shape under the projection?}
\lecturerresponse{There are two different statements. One is that circles project to circles. The other is the weaker statement that sufficiently small shapes preserve their local shape because angles are preserved. The mapping is conformal: it preserves angles, though not relative sizes.}
\end{classroomqa}

\section{Uniform Negative Curvature}

The uniformly negatively curved case is introduced through its metric.

The negatively curved analog of the sphere is hyperbolic space. In the same recursive spirit as before, its metric is obtained by replacing \(\sin\alpha\) with \(\sinh\alpha\):
\begin{equation}
d\mathcal H_n^2=d\alpha^2+\sinh^2\alpha\,d\Omega_{n-1}^2.
\end{equation}
The lower-dimensional angular piece remains the unit sphere metric \(d\Omega_{n-1}^2\).

\begin{classroomqa}
\lecturetimestamp{00:55:51}
\audiencequestion{In the metric, is that lower-dimensional differential still the spherical one?}
\lecturerresponse{Yes. That lower-dimensional piece is the ordinary unit-sphere metric. The change is in the \(\sinh^2\alpha\) factor in front of it.}
\end{classroomqa}

The hyperbolic sine is
\begin{equation}
\sinh\alpha=\frac{e^\alpha-e^{-\alpha}}{2}.
\end{equation}
The minus sign matters; in the lecture it is corrected immediately by a student.

There is also one crucial difference in the range of the coordinate. For the sphere, the new coordinate runs from one end to the other and stops at \(\pi\). Here,
\begin{equation}
0\leq \alpha<\infty.
\end{equation}
The ordinary sine rises and then comes back to zero. The hyperbolic sine begins the same way near \(\alpha=0\), but then it grows exponentially.

That means the lower-dimensional spheres around us grow bigger and bigger as we move outward. In this geometry the sky at increasing radius does not first grow and then shrink. It keeps growing.

In three spatial dimensions the corresponding FRW form is
\begin{equation}
d\tau^2=dt^2-a(t)^2\,d\mathcal H_3^2,
\end{equation}
and the sign label is
\begin{equation}
k=-1.
\end{equation}

\section{The Disk Picture and What It Means}

Negative curvature becomes intuitive in the disk picture made famous by Escher, especially \emph{Limit Circle Number Four}. The analogy with stereographic projection is useful, but the behavior is opposite.

For the sphere, the projection looks infinite while the real geometry is finite. For the negatively curved case, the disk picture looks finite while the real geometry is infinite. Equal ``continents'' or ``galaxies'' appear smaller and smaller as one moves toward the boundary, but geometrically they are all the same size. If distance is measured by counting how many equal pieces one crosses, then the boundary is infinitely far away.

Thus a known object in negative curvature looks too small relative to flat space. It subtends too small an angle.

The geometry is still homogeneous. Although the picture makes the outer pieces look tiny, any observer moved to another location would see the same intrinsic geometry.

\begin{classroomqa}
\lecturetimestamp{01:06:42}
\audiencequestion{Would an observer in any one of those circles see a similar pattern?}
\lecturerresponse{Yes, exactly the same thing. The space is homogeneous.}
\end{classroomqa}

\begin{classroomqa}
\lecturetimestamp{01:07:03}
\audiencequestion{In the spherical case distant things seem too big, while here they seem too small?}
\lecturerresponse{Yes. In positive curvature a known object subtends too large an angle relative to flat space. In negative curvature it subtends too small an angle.}
\end{classroomqa}

\begin{classroomqa}
\lecturetimestamp{01:08:09}
\audiencequestion{Is the relevance of the stereographic-style projection that an astronomer can reconstruct it from observations?}
\lecturerresponse{Yes. The astronomer is effectively a surveyor who surveys triangles in the sky. For the immediate geometric point one can set aside the time dependence and focus on angles and distances.}
\end{classroomqa}

There is also a useful local picture. A small patch of negative curvature can be drawn as a saddle surface in ordinary three-dimensional space. But only small patches can be drawn that way; the whole hyperbolic plane cannot be embedded globally as a nice ordinary surface in three-dimensional Euclidean space.

\section{The Flat Case}

Between the positively curved and negatively curved cases sits flat space. One may think of it as a limit, but it is enough simply to think of ordinary flat Euclidean space.

Its FRW form is
\begin{equation}
d\tau^2=dt^2-a(t)^2\left(dx^2+dy^2+dz^2\right),
\end{equation}
with
\begin{equation}
k=0.
\end{equation}

These are the three homogeneous FRW geometries:
\begin{align}
k=+1 &: \text{positive spatial curvature},\\
k=0 &: \text{flat space},\\
k=-1 &: \text{negative spatial curvature}.
\end{align}

\begin{classroomqa}
\lecturetimestamp{01:15:07}
\audiencequestion{If one projects the geometry onto a plane and follows a signal, would the projected speed of light vary from place to place?}
\lecturerresponse{Yes, on the projection it would. On the spherical projection it would appear to go faster and faster farther out; on the negative-curvature disk it would appear to go slower and slower near the edge. That variation belongs to the picture, not to the underlying physics.}
\end{classroomqa}

\begin{classroomqa}
\lecturetimestamp{01:17:05}
\audiencequestion{Has one actually used looking at far-away objects in this way to determine flatness?}
\lecturerresponse{Yes, but it is really the angular information that is used. One does not determine flatness by trying to measure how fast light appears to move on the projection.}
\end{classroomqa}

\section{Back to the Cosmological Equations}

After the break comes the connection between the cosmological equations already studied and Einstein's equations. At this point the board still carries the flat spatial metric fragment.

\begin{figure}[htbp]
\centering
\includegraphics[width=0.72\textwidth]{lecture_05_figure_02.png}
\caption{Flat FRW metric fragment.}
\lectureframe{01:18:18}
\end{figure}

A very brief Newtonian review comes first. Put us at the center. Look at a galaxy moving away according to the Hubble law. By Newton's theorem for a spherically symmetric distribution, only the mass inside the sphere whose radius matches the galaxy's distance contributes to the force on it. Then there are three possibilities: the galaxy moves above the escape velocity, below the escape velocity, or exactly at the escape velocity.

What is recalled next is that for ordinary nonrelativistic matter the density goes like the mass divided by \(a^3\), for radiation it goes like \(1/a^4\), and for vacuum energy it does not change as the sphere changes.

In the earlier Newtonian treatment, the cosmological equation had a density term and, if one did not sit exactly at the escape-velocity case, an additional constant term proportional to \(1/a^2\). The earlier discussion mostly set that term to zero. Empirically it is small, but conceptually there is no reason to leave it out.

\section{Einstein's Equations and the \(00\) Component}

The next step is Einstein's field equations. On the left-hand side stands the geometry, built from the spacetime curvature. On the right-hand side stands the energy-momentum tensor, the source of the gravitational field. In the lecture's normalization one writes
\begin{equation}
R^{\mu\nu}-\frac{1}{2}g^{\mu\nu}R=4\pi G\,T^{\mu\nu}.
\end{equation}
\editorialnote{In the visible board line the factor \(\tfrac12\) is not legible, but around 01:33:36 the spoken correction restores it. The equation here keeps that correction while preserving the lecture's \(4\pi G\) normalization.}

Not all ten independent components are needed here. By this stage so much has already been assumed --- homogeneity, the FRW form, and a single time-dependent scale factor --- that there is only one arbitrary function left, namely \(a(t)\). The natural component to choose is the time-time component, because \(T^{00}\) is the energy density \(\rho\).

\begin{figure}[htbp]
\centering
\includegraphics[width=0.62\textwidth]{lecture_05_figure_05.png}
\caption{Einstein field equation and its \(00\) component.}
\lectureframe{01:30:54}
\end{figure}

So the focus is
\begin{equation}
R^{00}-\frac{1}{2}g^{00}R=4\pi G\,\rho.
\end{equation}

The right-hand side is the energy density. The left-hand side is geometry.

\begin{classroomqa}
\lecturetimestamp{01:39:09}
\audiencequestion{What about the other components of the energy-momentum tensor?}
\lecturerresponse{They do not vanish in general. The space-space components can represent pressure. In the matter-dominated case the pressure is zero, in the radiation-dominated case it is one-third of the energy density, and for vacuum energy one writes \(P=w\rho\) with \(w=-1\).}
\end{classroomqa}

\begin{classroomqa}
\lecturetimestamp{01:40:24}
\audiencequestion{So if one had only matter, would the other components of \(T^{\mu\nu}\) vanish?}
\lecturerresponse{For pressureless matter, yes, the pressure components vanish.}
\end{classroomqa}

\section{What the Einstein Calculation Produces}

The full curvature calculation is not carried out on the board. What matters is what appears when an FRW metric is substituted into the Einstein tensor.

One kind of term comes just from the fact that the scale factor depends on time. Even if space were flat, a time-dependent scale factor would still give spacetime curvature. That contribution is proportional to
\begin{equation}
\left(\frac{\dot a}{a}\right)^2.
\end{equation}

The second kind of term comes from the spatial curvature itself. Even if the scale factor were frozen in time, curved space would still contribute to the spacetime curvature. For a simple uniformly curved geometry, curvature goes like one over the radius of curvature squared, so it contributes a term proportional to
\begin{equation}
\frac{k}{a^2},
\end{equation}
with a positive sign for the sphere, a negative sign for the hyperbolic case, and zero for flat space.

\begin{classroomqa}
\lecturetimestamp{01:36:53}
\audiencequestion{Should there also be a factor such as a radius of curvature?}
\lecturerresponse{That is already the \(1/a^2\) factor. For these simple uniformly curved geometries, curvature is just the inverse radius of curvature squared.}
\end{classroomqa}

The coefficient bookkeeping is done on the fly and is partly checked against the earlier Newtonian answer. The endpoint is\footnote{This displayed Friedmann form consolidates the spoken, self-corrected coefficient discussion and the comparison with the earlier Newtonian result around 01:32:35--01:36:53.}
\begin{equation}
\left(\frac{\dot a}{a}\right)^2+\frac{k}{a^2}=\frac{8\pi G}{3}\rho,
\end{equation}
or equivalently
\begin{equation}
\left(\frac{\dot a}{a}\right)^2=\frac{8\pi G}{3}\rho-\frac{k}{a^2}.
\end{equation}

This is the form that links the three FRW geometries to the earlier cosmological dynamics.

\begin{classroomqa}
\lecturetimestamp{01:37:52}
\audiencequestion{Could \(k\) take any value if one redefined \(a\)?}
\lecturerresponse{In the flat case there is no unique natural scale for \(a\), so one can rescale it, but \(k\) remains zero. In the spherical and negatively curved cases there is a natural curvature scale. If one takes \(a(t)\) to be that natural radius of curvature, then \(k\) is just \(+1\) or \(-1\).}
\end{classroomqa}

\begin{classroomqa}
\lecturetimestamp{01:41:00}
\audiencequestion{Is \(k\) equal to \(+1\), \(-1\), or \(0\) in that final equation?}
\lecturerresponse{Yes. It is \(+1\) for the sphere, \(-1\) for the hyperbolic case, and \(0\) for flat space.}
\end{classroomqa}

Newton's treatment did not lead us astray; it gave the same basic equation. What general relativity adds is the interpretation. From the Newtonian point of view, the extra constant distinguishes above escape velocity, below escape velocity, and the knife-edge case. From the Einstein point of view, the same sign tells us the geometry of space itself.

\begin{classroomqa}
\lecturetimestamp{01:42:46}
\audiencequestion{How did this work before one included a cosmological constant?}
\lecturerresponse{Before the cosmological constant, there was a tight relationship between escape-velocity behavior and spatial geometry. If \(k=+1\), space is positively curved and the motion is below escape velocity, so the universe can expand and then recontract. If \(k=-1\), space is negatively curved and the motion is above escape velocity, so the universe keeps expanding. If \(k=0\), one has the marginal knife-edge case.}
\end{classroomqa}

The same sign \(k\) can therefore be read in two ways: as a dynamical classification of the motion and as the sign of the curvature of space.
